\subsection{Implementation Roadmap and Challenges}
\label{sec:implementation-roadmap}

\subsubsection{Phased deployment strategy}

The personalization system is designed for incremental deployment:

\paragraph{Phase 1: Data collection infrastructure (Week 1--2)}
\begin{itemize}
    \item Extend database schema with new tables;
    \item Implement asynchronous behavior logging;
    \item Add UserContext data model and database mapping;
    \item Begin collecting baseline user interaction data.
\end{itemize}

\paragraph{Phase 2: Offline model training (Week 3--4)}
\begin{itemize}
    \item Implement feature extraction pipeline;
    \item Train initial per-user models using collected data;
    \item Evaluate model performance on test sets;
    \item Establish baseline metrics for comparison.
\end{itemize}

\paragraph{Phase 3: Integration with HERA (Week 5--6)}
\begin{itemize}
    \item Create PersonalizationAgent class;
    \item Implement model loading and inference logic;
    \item Integrate UserContext injection into orchestrator;
    \item Deploy passive prediction (no suggestions yet, only logging).
\end{itemize}

\paragraph{Phase 4: Proactive suggestions (Week 7--8)}
\begin{itemize}
    \item Enable background inference and suggestion generation;
    \item Implement suggestion feedback collection;
    \item Add conflict resolution logic;
    \item Monitor acceptance rates and tune thresholds.
\end{itemize}

\subsubsection{Technical challenges and mitigations}

\paragraph{Cold start problem}
New users have no historical data. \emph{Mitigation:} Use a global
default model trained on all users, then transition to personal model
after collecting 50+ interactions. Explicitly communicate to new users
that personalization improves over time.

\paragraph{Concept drift}
User preferences may change (e.g., seasonal temperature preferences,
lifestyle changes). \emph{Mitigation:} Implement exponential decay
weighting where recent data has higher influence. Monitor validation
performance and retrain when drift is detected.

\paragraph{Privacy concerns}
Detailed behavior logging raises privacy considerations. \emph{Mitigation:}
Store data locally (not cloud), implement data retention policies
(auto-delete logs older than 6 months), provide user dashboard to review
and delete their data, clearly communicate what is being logged.

\paragraph{Model debugging}
When predictions are wrong, debugging is harder than rule-based systems.
\emph{Mitigation:} Log model predictions alongside actual user actions,
implement feature importance analysis, provide admin dashboard showing
learned patterns and model confidence scores.

\subsection{Comparison with Original TinyML Approach}

Table~\ref{tab:old-vs-new} contrasts the original TinyML thermal comfort
prediction with the new multi-user personalization pipeline.

\begin{table}[H]
    \centering
    \caption{Comparison of original and redesigned ML approach.}
    \label{tab:old-vs-new}
    \renewcommand{\arraystretch}{1.2}
    \small
    \begin{tabularx}{\textwidth}{@{} >{\raggedright\arraybackslash}p{3cm} X X @{}}
        \toprule
        \textbf{Aspect} & \textbf{Original: TinyML Thermal Comfort} & \textbf{New: Multi-User Personalization} \\
        \midrule
        Deployment target &
        ESP32 microcontroller (on-device inference) &
        Backend server (Python runtime) \\
        
        Prediction task &
        3-class comfort classification (cooler/comfortable/warmer) &
        Multi-task: next action prediction + preferred device states \\
        
        User modeling &
        Single shared model for all users &
        Per-user models with individualized patterns \\
        
        Data source &
        ASHRAE public dataset + optional user feedback &
        Implicit learning from all user actions (no manual labeling) \\
        
        Model architecture &
        Tiny MLP (8 hidden units, TFLite quantized) &
        LightGBM (100+ trees, full precision) \\
        
        Integration &
        MQTT telemetry → model → suggestion &
        Deep integration with HERA agents via UserContext \\
        
        Personalization scope &
        Limited to thermal comfort preferences &
        All device control patterns, timing, and preferences \\
        
        Actionability &
        Reactive: only responds to current sensor readings &
        Proactive: suggests actions before user requests \\
        
        Memory constraints &
        Strict ($<$100KB model size) &
        Relaxed (multiple MB acceptable) \\
        
        Training frequency &
        Offline, manual retraining &
        Automated nightly batch retraining \\
        \bottomrule
    \end{tabularx}
\end{table}

The redesigned approach shifts focus from resource-constrained edge
inference to rich, personalized behavior modeling. This better aligns with
the actual system capabilities (backend has sufficient compute resources)
and user needs (personalization across all interactions, not just thermal
comfort).

\subsection{Related Work and Academic Context}

The multi-user personalization pipeline draws inspiration from several
research domains:

\subsubsection{Smart home behavior learning}
Prior work in smart home research has explored learning user routines for
automation. Notable examples include:
\begin{itemize}
    \item \textbf{Activity recognition:} Using sensors to detect user
    activities and trigger context-aware actions \cite{cook_smart_homes_2012}.
    However, these systems typically require dense sensor deployments.
    Our approach learns from explicit device commands rather than passive
    sensing.
    \item \textbf{Routine mining:} Discovering temporal patterns in home
    automation logs \cite{roy_routine_mining_2016}. We extend this by
    combining temporal patterns with environmental context and per-user
    modeling.
\end{itemize}

\subsubsection{Recommender systems}
The per-user model architecture parallels collaborative filtering
approaches in recommender systems:
\begin{itemize}
    \item \textbf{Netflix-style embeddings:} For large-scale systems
    ($>$1000 users), a shared model with user embeddings is more efficient
    \cite{netflix_embeddings_2016}. Our family-scale deployment
    (3--5 users) benefits more from dedicated per-user models.
    \item \textbf{Context-aware recommendation:} Incorporating time-of-day,
    environmental context, and sequential patterns into recommendations
    \cite{context_aware_rec_2015}.
\end{itemize}

\subsubsection{Multi-agent personalization}
The integration of ML personalization with LLM-based multi-agent systems
is a relatively novel contribution:
\begin{itemize}
    \item Most multi-agent frameworks focus on task decomposition, not
    user-specific adaptation \cite{langchain_agents_2023}.
    \item Our UserContext injection approach enables personalization
    without requiring separate agent instances per user, balancing
    flexibility with resource efficiency.
\end{itemize}

\subsection{Summary and Contributions}

This chapter introduced a comprehensive multi-user behavior learning and
personalization pipeline for the HERA smart home system. The key
contributions are:

\begin{enumerate}
    \item \textbf{Clear separation of concerns:} Distinguishing
    safety-critical rule-based logic from learned personalization,
    resolving the architectural confusion of the original TinyML anomaly
    detector.
    
    \item \textbf{Per-user modeling at family scale:} Demonstrating that
    individual models outperform shared models for small user populations,
    with practical implementation via lazy loading and LRU caching.
    
    \item \textbf{Hybrid reactive-proactive system:} Combining reactive
    personalization (applying learned preferences to user commands) with
    proactive suggestions (predicting user needs before they ask).
    
    \item \textbf{LLM-ML integration:} Showing how learned behavioral
    models can enhance LLM-based agents through UserContext injection,
    enabling personalization without fragmenting the agent architecture.
    
    \item \textbf{Multi-user conflict resolution:} Addressing the practical
    challenge of shared device control in multi-user households through
    role-based priorities and learned preference reconciliation.
    
    \item \textbf{Privacy-preserving design:} Keeping all behavior data
    local, implementing data retention policies, and providing user
    transparency through learned pattern dashboards.
\end{enumerate}

Unlike the original TinyML thermal comfort predictor—which addressed a
narrow use case with a resource-constrained deployment target—the
personalization pipeline provides broad value across all user
interactions. It transforms HERA from a command executor into an adaptive
assistant that learns each family member's unique preferences and proactively
supports their daily routines.

The implementation roadmap provides a practical path to deployment,
starting with data collection infrastructure and incrementally adding
model training, integration, and proactive features. This phased approach
allows validation of each component before proceeding, reducing deployment
risk while building toward a comprehensive personalized smart home
experience.
